% Loki III as one flow, from input to score. Stages are boxes, features are listed under their stage, the data
% on each arrow is written on it, and each numbered mark is a main-text argument (the circled text is the
% number of its subsection) drawn on the stage that it tests. Shared palette: IMAGE/odin3-tikz-styles.tex.
\newcommand{\lkmark}[3]{\node[circle, draw=o3grn8, fill=white, line width=0.5pt, inner sep=0.6pt,
    font=\tiny, text=o3ink] at (#1) {\ref{#2}#3};}
\begin{tikzpicture}[odin, x=1cm, y=1cm,
  stage/.style={box, light, text width=29mm, minimum height=14mm, font=\footnotesize\bfseries},
  feat/.style={plain, text width=29mm, font=\scriptsize, color=o3sub, align=left, anchor=north},
  closed/.style={box, external, text width=29mm, font=\scriptsize}]
  % input and answer key
  \node[box, neutral, text width=27mm, minimum height=14mm] (odin) at (0,0)
    {\textbf{Odin III}\\\odinsub{netlist and provenance answer key}};
  \node[feat, text width=27mm] at (0,-0.9) {HDL source, floors f1 to f3 (LUT, hard blocks, gates)};
  % stages
  \node[stage] (l0) at (3.9,0) {L0 datasets};
  \node[feat] at (3.9,-0.9) {naming-neutral ids; renamed twins; labels from the key; 68 micro and 156 VerilogEval designs};
  \node[stage] (l1) at (7.6,0) {L1 Bayesian flow};
  \node[feat] at (7.6,-0.9) {detectors, provers, calibration; hard blocks, adders, muxes, registers};
  \node[stage] (l2) at (11.3,0) {L2 graph analysis};
  \node[feat] at (11.3,-0.9) {clusters, GNN against 1-WL; loss audit with the provenance map};
  \node[stage] (l3) at (15.0,0) {L3 blocks};
  \node[feat] at (15.0,-0.9) {anchors, grow to path end, join, Gaussian kind model};
  \node[box, neutral, text width=27mm, minimum height=14mm] (out) at (18.6,0)
    {\textbf{Recovered blocks and kinds}\\\odinsub{scored against the key}};
  \node[closed] (wf) at (11.3,2.6) {\textbf{closed:} word-first detector, split step, clusters as hierarchy};
  % flow arrows with the data on them
  \draw[flow] (odin) -- node[edgelabel, above=1pt] {netlist} (l0);
  \draw[flow] (l0) -- node[edgelabel, above=1pt] {graph.json} (l1);
  \draw[flow] (l1) -- node[edgelabel, above=1pt] {claims} (l2);
  \draw[flow] (l2) -- node[edgelabel, above=1pt] {signatures} (l3);
  \draw[flow] (l3) -- node[edgelabel, above=1pt] {blocks} (out);
  \draw[ref] (l2.north) -- (wf.south);
  % the answer key runs under the whole pipeline to the score
  \draw[ref] (odin.west) -- (-1.9,0) -- (-1.9,-3.3) -- (18.6,-3.3) -- (out.south);
  \node[edgelabel] at (2.2,-3.3) {answer key};
  % LLM side path: it takes the netlist and, in arm B, the claims of L1 as hints
  \node[box, external, text width=33mm, minimum height=12mm] (gb) at (9.0,-5.4)
    {\textbf{Gate-Breaker harness}\\\odinsub{arms A, B, C; three draws}};
  \node[box, external, text width=30mm, minimum height=12mm] (rtl) at (14.2,-5.4)
    {\textbf{Model RTL}\\\odinsub{pass at k against the reference}};
  \fill (5.75,0) circle (1pt); \fill (9.45,0) circle (1pt);
  \draw[flow] (5.75,0) -- (5.75,-5.4) -- node[edgelabel, pos=0.3, anchor=west] {gate netlist} (gb.west);
  \draw[flow] (9.45,0) -- (9.45,-3.6) -| (gb.north);
  \node[edgelabel, anchor=west] at (9.5,-3.9) {claims as hints, arm B only};
  \draw[flow] (gb) -- (rtl);
  % argument marks (circled: subsection number in the Results and in the language-model sections)
  \lkmark{3.0,0.9}{arg:ra1}{}
  \lkmark{4.8,0.9}{arg:ra5}{}
  \lkmark{6.7,0.9}{arg:ra2}{}
  \lkmark{7.5,0.9}{arg:ra3}{}
  \lkmark{10.4,0.9}{arg:ra4}{}
  \lkmark{11.3,0.9}{arg:ra6}{}
  \lkmark{12.2,0.9}{arg:ra7}{}
  \lkmark{14.2,0.9}{arg:ra8}{}
  \lkmark{14.7,0.9}{arg:ra9}{}
  \lkmark{15.2,0.9}{arg:ra10}{}
  \lkmark{15.7,0.9}{arg:ra11}{}
  \lkmark{16.2,0.9}{arg:ra12}{}
  \lkmark{16.7,0.9}{arg:ra13}{}
  \lkmark{7.8,-4.8}{arg:la1}{}
  \lkmark{8.6,-4.8}{arg:la2}{}
  \lkmark{9.4,-4.8}{arg:la3}{}
  \lkmark{10.2,-4.8}{arg:la4}{}
\end{tikzpicture}
